
Static analysis metrics, originally designed for human-written code, correlate with functional correctness of LLM-generated code. This study quantifies that relationship, finding that pylint score alone achieves r=0.873 correlation with pass@1 on HumanEval/MBPP, a result that challenges assumptions about the necessity of multi-metric ensemble approaches.

The deployment of LLM-generated code presents a practical problem: running test suites on k candidate completions costs $k\times$ compute, yet accepting unchecked code risks deploying incorrect implementations. Static analysis tools such as pylint, mypy, and radon are fast and deterministic, but prior work has used them primarily for feedback and correction rather than prediction. If static analysis metrics could predict functional correctness, practitioners could filter LLM outputs at substantially reduced cost.

Recent work demonstrates that iterative static analysis feedback improves code quality. Blyth et al. (2025) show that Bandit and Pylint feedback reduces security issues from over 40\% to 13\% on HumanEval/MBPP. CodeQUEST \citep{Liu2025CodeQUEST} reports ''meaningful correlation'' between static analysis metrics and LLM code quality. However, no prior study quantifies the predictive correlation with regression coefficients or p-values. Studies use static analysis for correction but not prediction.

This work addresses the correlation gap directly. The central hypothesis is that static analysis quality metrics are predictive of LLM-generated code correctness. The experiments test three research questions:

\begin{enumerate}
\item Does pylint score achieve $r\geq 0.35$ correlation with pass@1 after controlling for code length?
\item Does weighted ensemble combination of static analysis metrics outperform individual metrics?
\item Does the correlation generalize across different LLMs with low variance?
\end{enumerate}

The contributions are:

\textbf{Empirical:} A quantified correlation study establishing r=0.873 between pylint scores and functional correctness for LLM-generated code on standard benchmarks.

\textbf{Methodological:} A partial correlation framework controlling for code length as a confounding variable. The minimal difference between raw (r=0.868) and partial (r=0.873) correlations indicates that the static analysis signal is not an artifact of code brevity.

\textbf{Practical:} Evidence that pylint alone is sufficient for quality filtering. Weighted ensemble combination degrades performance (r=0.861 < r=0.873), simplifying deployment to a single metric.
